%! Unit Opener: Functions as Models — Unit 2, Lesson 2.0 — Homework (Answer Key)
\documentclass[10pt]{article}
\usepackage{atda-article}
\usepackage{atda-key}   % pulls in -boxes; do NOT also load -boxes
\newcommand{\TallMath}[1]{$\displaystyle #1\rule[-1.4em]{0pt}{3.2em}$}
\setlength{\workrowsep}{8pt}

\begin{document}

\pageheader{Unit 2, Lesson 2.0}{Homework --- Answer Key}
% NAMESTRIP: no \namedateperiod here — mirrors the blank. Teacher-only prose goes
% in the lesson plan as \begin{teachernote}[Homework], never in a key.

\vspace{0.15in}

\boxguard[24]
\begin{notesbox}{Practice --- The Skills Unit 2 Is Built On}
\small
Items 1--4 are the diagnostic skills again with new numbers. Show your work.

\begin{enumerate}[leftmargin=1.6em, itemsep=6pt, topsep=4pt]

\item Let $g(x)=5x-8$. Find $g(-2)$, then find the value of $x$ for which $g(x)=7$.
\begin{work}
  g(-2) &= 5(-2)-8 \\
        &= -18 \\
  5x-8  &= 7 \\
    5x  &= 15 \\
     x  &= 3
\end{work}

\item For which value of $x$ is $\dfrac{x-6}{x+2}$ undefined?
\begin{work}
  x+2 &= 0 \\
    x &= -2
\end{work}

\item Find the $x$-intercept and the $y$-intercept of $4x-5y=20$. Write each as an ordered pair.
\begin{work}
  \text{set } y=0: \quad 4x &= 20 \\
                         x &= 5 \\
  \text{set } x=0: \quad -5y &= 20 \\
                          y &= -4
\end{work}

\item A tank is draining. Its volume $V$ (gallons) after $t$ minutes:
\smallskip

\renewcommand{\arraystretch}{1.25}
\begin{tabular}{@{}|c|c|c|c|c|@{}}
\hline
$t$ (minutes) & $0$ & $2$ & $4$ & $6$ \\ \hline
$V$ (gallons) & $250$ & $220$ & $190$ & $160$ \\ \hline
\end{tabular}
\smallskip

By how much does $V$ change per minute, and what is $V$ when $t=10$?
\begin{work}
  \text{change per minute} &= \dfrac{220-250}{2} \\
                           &= -15 \text{ gallons per minute} \\
                     V(10) &= 250-15(10) \\
                           &= 100 \text{ gallons}
\end{work}

\end{enumerate}
\end{notesbox}

\vspace{0.10in}

\boxguard[30]
\begin{scenariobox}[In Context --- A Video Takes Off]{royal}
\small
A club posts a video. The graph shows $V$, its total views \textbf{in hundreds}, $d$ days after
posting.

\begin{center}
\begin{tikzpicture}
\begin{axis}[
    width=9.8cm, height=4.8cm,
    xlabel={$d$ (days after posting)}, ylabel={$V$ (views, hundreds)},
    xmin=0, xmax=4.4, ymin=0, ymax=18,
    xtick={0,1,2,3,4}, ytick={0,4,8,12,16},
    grid=both, grid style={linegray!45}, axis lines=left,
    tick label style={font=\scriptsize}, label style={font=\scriptsize}]
  \addplot[royal, very thick, domain=0:4, samples=60]{2^x};
\end{axis}
\end{tikzpicture}
\end{center}

\begin{enumerate}[leftmargin=1.6em, itemsep=6pt, topsep=2pt, start=5]

\item Where does the graph meet the vertical axis, and how many actual views is that?

\ansline{At $(0,1)$ --- one hundred views on the day it was posted.}\\

\item Is $V$ increasing or decreasing on $0<d<4$? Does this graph have an absolute maximum over
the four days shown? Explain.

\ansline{Increasing the whole way. Its absolute maximum over these four days is}\\
\ansline{$V=16$ (1600 views) at $d=4$, the right end --- it is still climbing there.}\\

\item Read $V$ off the graph at $d=0,1,2,3,4$ and fill in the table. Then say which family this
is, and what in the table proves it is \emph{not} linear.
\smallskip

\renewcommand{\arraystretch}{1.4}
\begin{tabular}{@{}|c|c|c|c|c|c|@{}}
\hline
$d$ & $0$ & $1$ & $2$ & $3$ & $4$ \\ \hline
$V$ & \ans{1} & \ans{2} & \ans{4} & \ans{8} & \ans{16} \\ \hline
\end{tabular}
\smallskip

\ansline{Exponential: each day the views are multiplied by $2$.}\\
\ansline{Not linear --- the amounts added are $1, 2, 4, 8$, never a constant.}\\

\item Estimate the day the video passes $1000$ views, and say how you used the graph.

\ansline{Between day 3 and day 4: $1000$ views is $V=10$, and the curve passes $10$ there.}\\
\end{enumerate}
\end{scenariobox}

\vspace{0.10in}

% [20] was tested and broke this box, stranding item 9(c) alone atop p3.
% The box measures ~23 baselines, so the guard must exceed that.
\boxguard[24]
\begin{scenariobox}[A First Look Ahead]{goldacc}
\small
\begin{enumerate}[leftmargin=1.6em, itemsep=6pt, topsep=2pt, start=9]

\item In Lesson 1.7 you found the excluded values of a rational expression. Let
$f(x)=\dfrac{x+1}{x-3}$.
\begin{enumerate}[label=(\alph*), leftmargin=1.6em, itemsep=5pt, topsep=4pt]
  \item Which input must be left out, and why? That input is missing from the
        \textbf{domain} of $f$ --- Lesson 2.1's word for it.
\begin{work}
  x-3 &= 0 \\
    x &= 3
\end{work}

  \item Here are four outputs near that input. What is happening to $f(x)$?
        \smallskip

        \renewcommand{\arraystretch}{1.25}
        \begin{tabular}{@{}|c|c|c|c|c|@{}}
        \hline
        $x$ & $2.9$ & $2.99$ & $3.01$ & $3.1$ \\ \hline
        $f(x)$ & $-39$ & $-399$ & $401$ & $41$ \\ \hline
        \end{tabular}
        \smallskip

\ansline{The outputs blow up --- huge negative just left of $3$, huge positive just right.}\\

  \item The graph of $f$ runs alongside the vertical line $x=3$ forever without ever touching it.
        Lesson 2.5 has a name for that line. Guess it, or describe it in your own words.

\ansline{A vertical asymptote (accept any description of the graph shooting off at $x=3$).}\\
\end{enumerate}

\end{enumerate}
\end{scenariobox}

\vspace{0.10in}

\boxguard[22]
\begin{extensionbox}
\small
\textbf{Two savings plans.} Plan A starts at \$100 and adds \$50 every month. Plan B starts at
\$20 and \emph{doubles} every month. Fill in the table, then answer.

\smallskip
\renewcommand{\arraystretch}{1.4}
\begin{tabular}{@{}|c|c|c|c|c|c|@{}}
\hline
Month $m$ & $0$ & $1$ & $2$ & $3$ & $4$ \\ \hline
Plan A & \ans{100} & \ans{150} & \ans{200} & \ans{250} & \ans{300} \\ \hline
Plan B & \ans{20} & \ans{40} & \ans{80} & \ans{160} & \ans{320} \\ \hline
\end{tabular}
\smallskip

Plan A is ahead for the first three months. In which month does Plan B pass it, and what does
that tell you about adding a constant versus multiplying by one?

\ansline{Month $4$: \$320 beats \$300. A constant multiplier eventually overtakes}\\
\ansline{any constant addition, no matter how far behind it starts.}\\
\end{extensionbox}

\vspace{0.10in}

\begin{remindbox}
\small
\textbf{Next: Lesson 2.1 --- Function Fundamentals: Notation, Domain, and Range in Context}
(\texttt{AFDA.AF.2a, e}). Today you read domain and range off a picture. Lesson 2.1 makes them
precise, gives you the notation to write them down, and shows how a situation --- or an excluded
value --- cuts a domain short.
\end{remindbox}

\end{document}
